\documentclass[12pt,letterpaper]{article}
\usepackage{graphicx,textcomp}
\usepackage{natbib}
\usepackage{setspace}
\usepackage{fullpage}
\usepackage{color}
\usepackage[reqno]{amsmath}
\usepackage{amsthm}
\usepackage{fancyvrb}
\usepackage{amssymb,enumerate}
\usepackage[all]{xy}
\usepackage{endnotes}
\usepackage{lscape}
\newtheorem{com}{Comment}
\usepackage{float}
\usepackage{hyperref}
\newtheorem{lem} {Lemma}
\newtheorem{prop}{Proposition}
\newtheorem{thm}{Theorem}
\newtheorem{defn}{Definition}
\newtheorem{cor}{Corollary}
\newtheorem{obs}{Observation}
\usepackage[compact]{titlesec}
\usepackage{dcolumn}
\usepackage{tikz}
\usetikzlibrary{arrows}
\usepackage{multirow}
\usepackage{xcolor}
\newcolumntype{.}{D{.}{.}{-1}}
\newcolumntype{d}[1]{D{.}{.}{#1}}
\definecolor{light-gray}{gray}{0.65}
\usepackage{url}
\usepackage{listings}
\usepackage{color}

\definecolor{codegreen}{rgb}{0,0.6,0}
\definecolor{codegray}{rgb}{0.5,0.5,0.5}
\definecolor{codepurple}{rgb}{0.58,0,0.82}
\definecolor{backcolour}{rgb}{0.95,0.95,0.92}

\lstdefinestyle{mystyle}{
	backgroundcolor=\color{backcolour},   
	commentstyle=\color{codegreen},
	keywordstyle=\color{magenta},
	numberstyle=\tiny\color{codegray},
	stringstyle=\color{codepurple},
	basicstyle=\footnotesize,
	breakatwhitespace=false,         
	breaklines=true,                 
	captionpos=b,                    
	keepspaces=true,                 
	numbers=left,                    
	numbersep=5pt,                  
	showspaces=false,                
	showstringspaces=false,
	showtabs=false,                  
	tabsize=2
}
\lstset{style=mystyle}
\newcommand{\Sref}[1]{Section~\ref{#1}}
\newtheorem{hyp}{Hypothesis}

\title{Problem Set 1}
\date{Due: October 8, 2026}
\author{Applied Stats/Quant Methods 1}

\begin{document}
	\maketitle
	
	\section*{Instructions}
	\begin{itemize}
	\item Please show your work! You may lose points by simply writing in the answer. If the problem requires you to execute commands in \texttt{R}, please include the code you used to get your answers. Please also include the \texttt{.R} file that contains your code. If you are not sure if work needs to be shown for a particular problem, please ask.
\item Your homework should be submitted electronically on GitHub.
\item This problem set is due before 23:59 on Thursday October 8, 2026. No late assignments will be accepted.
	\end{itemize}
	
	\vspace{1cm}
	\section*{Question 1: Education}

A school counselor was curious about the average of IQ of the students in her school and took a random sample of 25 students' IQ scores. The following is the data set:\\
\vspace{.5cm}

\lstinputlisting[language=R, firstline=39, lastline=39, breaklines=true]{PS01_EL.R}


\vspace{1cm}

\begin{enumerate}
	\item Find a 90\% confidence interval for the average student IQ in the school (As a hint, you first need the mean and SD to create a CI):\\
	\section{Answer For Question 1.1}
	
The first determining length which is the number of values that we have in this data set. 

	\lstinputlisting[language=R, firstline=44, lastline=44]{PSO1_EL.R} 

\textit{ANSWER FOR THE LENGTH OF Y IS: 25}

The second is to determine the mean of sample the scores. 

	\lstinputlisting[language=R, firstline=50, lastline=52]{PSO1_EL.R} 
 
 \textit{ANSWER OF MEAN SCORE FOR THE SCORES:98.44}
 
 The third is to determine the standard deviation of the scores.
 
	\lstinputlisting[language=R, firstline=59, lastline=61]{PSO1_EL.R}

 \textit{ANSWER OF SD SCORE FOR THE SCORES: 13.09287}
 
	I have calculated the  length, mean and SD. In order to determine the 
	confidence interval, the standard error will also have to be calculated.
 
 	\lstinputlisting[language=R, firstline=68, lastline=70]{PSO1_EL.R}
 	
\textit{ANSWER FOR THE STANDARD ERROR SCORES: 2.618575 }
 
Now having together the length, mean, standard deviation, and the standard error we can now calculate the confidence interval at 90%
 
 CI = mean +/- critical value * SE
 
The confidence interval is set at 90%
 
 Determining the qnorm determines the proportions (with NORMAL distributions) of the p-value 
 
	\lstinputlisting[language=R, firstline=84, lastline=85]{PSO1_EL.R}
	
	\textit{ANSWER TO THE UPPER QNORM IS:  1.644854}
	\textit{ANSWER TO THE LOWER QNORM IS: -1.644854}
 
 Calculating the CI
 
 	\lstinputlisting[language=R, firstline=92, lastline=98]{PSO1_EL.R}
 	\lstinputlisting[language=R, firstline=99, lastline=99]{PSO1_EL.R}
 	\textit{ANSWER FOR THE LOWER TAIL: 94.13283}
 	\lstinputlisting[language=R, firstline=101, lastline=101]{PSO1_EL.R}
 	\textit{ANSWER FOR THE UPPER TAIL: 102.7472}
 	
 	Final Answer for Question 1.1: At the condidance level of 90%
 	the upper score average is approximetley 103. While the lower average 
 	of the IQ score is 84.
 
	\item Next, the school counselor was curious  whether  the average student IQ in her school is higher than the average IQ score (100) among all the schools in the country.\\ 
	
	\noindent Using the same sample, conduct the appropriate hypothesis test with $\alpha=0.05$.
	
\end{enumerate}

\newpage

	\section*{Question 2: Political Economy}

\noindent Researchers are curious about what affects the amount of money communities spend on addressing homelessness. The following variables constitute our data set about social welfare expenditures in the USA. \\
\vspace{.5cm}


\begin{tabular}{r|l}
	\texttt{State} &\emph{50 states in US} \\
	\texttt{Y} & \emph{per capita expenditure on shelters/housing assistance in state}\\
	\texttt{X1} &\emph{per capita personal income in state} \\
	\texttt{X2} &  \emph{Number of residents per 100,000 that are "financially insecure" in state}\\
	\texttt{X3} &  \emph{Number of people per thousand residing in urban areas in state} \\
	\texttt{Region} &  \emph{1=Northeast, 2= North Central, 3= South, 4=West} \\
\end{tabular}

\vspace{.5cm}
\noindent Explore the \texttt{expenditure} data set and import data into \texttt{R}.
\vspace{.5cm}
\lstinputlisting[language=R, firstline=42, lastline=42]{PS01.R}  
\vspace{.5cm}
\begin{itemize}

\item
Please plot the relationships among \emph{Y}, \emph{X1}, \emph{X2}, and \emph{X3}? What are the correlations among them (you just need to describe the graph and the relationships among them)?
\vspace{.5cm}
\item
Please plot the relationship between \emph{Y} and \emph{Region}? On average, which region has the highest per capita expenditure on housing assistance?
\vspace{.5cm}
\item
Please plot the relationship between \emph{Y} and \emph{X1}? Describe this graph and the relationship. Reproduce the above graph including one more variable \emph{Region} and display different regions with different types of symbols and colors.
\end{itemize}


\end{document}
